\section{Candy}
\label{sec:candy}

\problemstatement{There are $n$ children standing in a line, and the $i$-th
child has a rating $\textit{ratings}[i]$. We must give each child at least
one candy, subject to the rule that any child with a strictly higher rating
than an \textbf{immediate neighbor} (the child directly to the left or
directly to the right) must receive strictly more candies than that
neighbor. We must find the \textbf{minimum total number of candies}
required to satisfy all of these constraints.}

\begin{patternbox}
\emph{``Each element must be strictly greater than its immediate neighbor
whenever its value is strictly greater''} -- a constraint defined purely in
terms of comparisons with the elements directly to the left and right -- is
the signature of a \emph{two-pass directional greedy}. Whenever a problem
states a rule of the form ``compare element $i$ only to $i-1$ and $i+1$, and
the rule is not symmetric (higher rating needs strictly more, but the rule
says nothing directly about how much more when ratings are equal or lower),''
the standard technique is to satisfy the left-neighbor rule and the
right-neighbor rule in two completely separate linear passes, then combine
the two results -- rather than trying to track both directions
simultaneously in a single pass, which is awkward because a child's
required candy count can be driven up by either side independently.
\end{patternbox}

\subsection*{Understanding the problem carefully}

The rule is purely \emph{local}: child $i$ needs more candies than child
$i-1$ if and only if $\textit{ratings}[i] > \textit{ratings}[i-1]$, and
child $i$ needs more candies than child $i+1$ if and only if
$\textit{ratings}[i] > \textit{ratings}[i+1]$. These are two independent
constraints, one looking left and one looking right, and a child's final
candy count must simultaneously satisfy both. This naturally suggests
solving the left-looking constraint and the right-looking constraint
separately, each with its own simple left-to-right or right-to-left greedy
sweep, and then merging the two answers by taking, for each child, whichever
requirement is stricter (i.e.\ larger).

\subsection*{The left-to-right pass}

\begin{ideabox}[Greedy Choice -- Left Pass]
Scan left to right. Give every child at least $1$ candy initially. If
$\textit{ratings}[i] > \textit{ratings}[i-1]$, this child must have strictly
more candy than the child to its left, so set $\textit{left}[i] \gets
\textit{left}[i-1] + 1$. Otherwise, the left-neighbor constraint places no
extra demand on child $i$ beyond the baseline of $1$ candy, so
$\textit{left}[i] \gets 1$.
\end{ideabox}

This single pass correctly satisfies every left-looking constraint, because
each child's requirement here only ever depends on the child immediately
before it, which has already been finalized by the time we reach the
current child -- there is no need to look further back, since any chain of
strictly increasing ratings to the left has already been correctly
propagated forward by the recursive nature of the update rule.

\subsection*{The right-to-left pass}

\begin{ideabox}[Greedy Choice -- Right Pass]
Symmetrically, scan right to left. Give every child at least $1$ candy
initially. If $\textit{ratings}[i] > \textit{ratings}[i+1]$, set
$\textit{right}[i] \gets \textit{right}[i+1] + 1$. Otherwise,
$\textit{right}[i] \gets 1$.
\end{ideabox}

\subsection*{Why taking the maximum of both passes is correct}

Once both arrays are computed, child $i$'s final candy count must be at
least $\textit{left}[i]$ (to satisfy the left-neighbor rule) and at least
$\textit{right}[i]$ (to satisfy the right-neighbor rule). Since these are
two independent lower bounds on the same quantity (the number of candies
given to child $i$), and we want the \emph{minimum} total, the smallest
value that satisfies both lower bounds simultaneously is exactly their
maximum:
\[
  \textit{candies}[i] = \max(\textit{left}[i], \textit{right}[i]).
\]
No smaller value could satisfy both constraints, and this value, taken for
every child simultaneously, does not interfere with any other child's
constraint, because each constraint here only relates two physically
adjacent children, and the relative order between $\textit{left}[i]$ and
$\textit{left}[i\pm1]$ (or $\textit{right}[i]$ and $\textit{right}[i\pm1]$)
already correctly encodes the required strict inequality wherever the
rating comparison demands it. Taking componentwise maxima of two arrays that
each individually satisfy their own one-sided constraint is guaranteed to
satisfy both constraints together, without ever needing to re-verify the
combined array.

\subsection*{Algorithm}

\begin{enumerate}
  \item Initialize arrays $\textit{left}[0..n-1]$ and
        $\textit{right}[0..n-1]$, each filled with $1$.
  \item Left pass: for $i \gets 1$ to $n-1$, if
        $\textit{ratings}[i] > \textit{ratings}[i-1]$, set
        $\textit{left}[i] \gets \textit{left}[i-1] + 1$.
  \item Right pass: for $i \gets n-2$ down to $0$, if
        $\textit{ratings}[i] > \textit{ratings}[i+1]$, set
        $\textit{right}[i] \gets \textit{right}[i+1] + 1$.
  \item Compute $\textit{total} \gets \sum_{i=0}^{n-1} \max(\textit{left}[i],
        \textit{right}[i])$.
  \item Return \textit{total}.
\end{enumerate}

\begin{algorithm}[H]
\caption{Candy Distribution}
\begin{algorithmic}[1]
\Procedure{Candy}{\textit{ratings}}
  \State $n \gets \text{length}(\textit{ratings})$
  \State $\textit{left}[0..n-1] \gets 1$, $\textit{right}[0..n-1] \gets 1$
  \For{$i \gets 1$ to $n-1$}
    \If{$\textit{ratings}[i] > \textit{ratings}[i-1]$}
      \State $\textit{left}[i] \gets \textit{left}[i-1] + 1$
    \EndIf
  \EndFor
  \For{$i \gets n-2$ \textbf{down to} $0$}
    \If{$\textit{ratings}[i] > \textit{ratings}[i+1]$}
      \State $\textit{right}[i] \gets \textit{right}[i+1] + 1$
    \EndIf
  \EndFor
  \State $\textit{total} \gets 0$
  \For{$i \gets 0$ to $n-1$}
    \State $\textit{total} \gets \textit{total} + \max(\textit{left}[i], \textit{right}[i])$
  \EndFor
  \State \Return \textit{total}
\EndProcedure
\end{algorithmic}
\end{algorithm}

\subsection*{Dry run}

\dryrun{Take $\textit{ratings} = [1, 0, 2]$.}

\textbf{Left pass.} $\textit{left} = [1, 1, 1]$ initially.
\begin{itemize}
  \item $i=1$: $\textit{ratings}[1]=0 > \textit{ratings}[0]=1$? No. Stays
        $\textit{left}[1]=1$.
  \item $i=2$: $\textit{ratings}[2]=2 > \textit{ratings}[1]=0$? Yes.
        $\textit{left}[2] = \textit{left}[1]+1 = 2$.
\end{itemize}
Result: $\textit{left} = [1, 1, 2]$.

\textbf{Right pass.} $\textit{right} = [1, 1, 1]$ initially.
\begin{itemize}
  \item $i=1$: $\textit{ratings}[1]=0 > \textit{ratings}[2]=2$? No. Stays
        $\textit{right}[1]=1$.
  \item $i=0$: $\textit{ratings}[0]=1 > \textit{ratings}[1]=0$? Yes.
        $\textit{right}[0] = \textit{right}[1]+1 = 2$.
\end{itemize}
Result: $\textit{right} = [2, 1, 1]$.

\textbf{Combine.} $\textit{candies} = [\max(1,2), \max(1,1), \max(2,1)] =
[2, 1, 2]$.

The total is $2+1+2 = 5$. This is correct: child $0$ (rating $1$) needs more
than child $1$ (rating $0$), so child $0$ gets $2$ and child $1$ gets the
baseline $1$; child $2$ (rating $2$) needs more than child $1$, so child $2$
gets $2$. No constraint links child $0$ and child $2$ directly (they are not
neighbors), so this is fully consistent and minimal.

\subsection*{C++ implementation}

\begin{lstlisting}
#include <bits/stdc++.h>
using namespace std;

int candy(vector<int>& ratings) {
    int n = ratings.size();
    vector<int> left(n, 1), right(n, 1);

    for (int i = 1; i < n; i++) {
        if (ratings[i] > ratings[i - 1]) {
            left[i] = left[i - 1] + 1;
        }
    }

    for (int i = n - 2; i >= 0; i--) {
        if (ratings[i] > ratings[i + 1]) {
            right[i] = right[i + 1] + 1;
        }
    }

    int total = 0;
    for (int i = 0; i < n; i++) {
        total += max(left[i], right[i]);
    }
    return total;
}

int main() {
    vector<int> ratings = {1, 0, 2};
    cout << "Minimum candies: " << candy(ratings) << endl;
    return 0;
}
\end{lstlisting}

\begin{complexitybox}
\textbf{Time Complexity.} We perform two linear passes (left-to-right and
right-to-left) plus one linear combination pass, each $O(n)$, giving an
overall time complexity of
\[
  O(n).
\]
\textbf{Space Complexity.} We use two auxiliary arrays, \textit{left} and
\textit{right}, each of size $n$, giving a space complexity of $O(n)$. This
can be reduced to $O(1)$ extra space with a more careful single-pass
peak-detection technique, though the two-array version is the clearest to
derive and verify correct.
\end{complexitybox}

\begin{takeawaybox}
When a constraint compares each element only to its immediate left and
right neighbors, and the comparison is directionally asymmetric (a strict
inequality matters only in one direction at a time), decompose the problem
into two independent one-directional sweeps -- one enforcing the
left-looking rule, one enforcing the right-looking rule -- and combine the
two results with a pointwise maximum (for a ``need at least this many''
constraint) or minimum (for a ``cannot exceed this many'' constraint).
\end{takeawaybox}
